% Part I: the complexity dichotomy of realizability (from the paperX note)

\section{The order-convex structure and its clone}

Realizability along $z$ has, as the cell $a$ varies, variable domains
the fibres $z^{-1}(a(v))$, which are \emph{order-convex}, and the single
binary constraint $\le$.  The natural template collecting all such
constraints is the following.

\begin{defn}
  A set $C\subseteq\B^m$ is \emph{order-convex} if $a,b\in C$ and $a\le
  c\le b$ imply $c\in C$.  Let $L_m$ be the relational structure on
  $\B^m$ with the order $\le$ and every order-convex subset as a unary
  relation, and $\Pol(L_m)$ its clone of polymorphisms.
\end{defn}

\begin{lem}\label{lem:polchar}
  $f\in\Pol(L_m)$ iff $f$ is monotone, idempotent, and
  \emph{order-convex-preserving}: $f(x_1,\dots,x_k)\in\conv\{x_1,\dots,
  x_k\}=\{c:\exists i,j\ x_i\le c\le x_j\}$.
\end{lem}
\begin{proof}
  Preserving $\le$ is monotonicity; singletons $\{a\}=[a,a]$ are
  order-convex, so preserving them is idempotence.  The convex hull is
  order-convex and contains the arguments, so preservation forces
  $f(\vec x)\in\conv\{\vec x\}$; conversely that condition makes $f$
  preserve every order-convex set containing all the $x_i$.
\end{proof}

An operation $w$ is a \emph{weak near-unanimity} (WNU) operation if it is
idempotent and $w(y,x,\dots,x)=w(x,y,x,\dots,x)=\cdots=w(x,\dots,x,y)$; we
write $m(x,y)=f(x,x,y)=f(x,y,x)=f(y,x,x)$ for a ternary one.  A finite
idempotent clone has \emph{bounded width} iff it has WNU polymorphisms of
all but finitely many arities, equivalently a ternary WNU $u$ and a
$4$-ary WNU $v$ with $u(x,x,y)=v(x,x,x,y)$ \cite{BartoKozik,KKVW}; it is
\emph{Taylor} (hence its CSP is in $\mathrm P$, by \cite{Bulatov,Zhuk})
iff it has a $4$-ary Siggers polymorphism $s(a,r,e,a)=s(r,a,r,e)$
\cite{Siggers,KMM}; otherwise the CSP is $\mathrm{NP}$-complete.  We
exclude a ternary WNU (Proposition~\ref{prop:base}) and a Siggers
operation (Proposition~\ref{prop:nosiggers}), so the conclusions below do
not depend on which equivalent formulation of bounded width one adopts.

\section{Reduction to the cube $\B^3$}

\begin{lem}[Order-convex restriction]\label{lem:restrict}
  Let $m\ge3$ and $T=\B^3\times\{0\}^{m-3}$, an interval hence
  order-convex.  Any $f\in\Pol(L_m)$ preserves $T$ and restricts to a
  polymorphism of $L_3$ on $T\cong\B^3$; the restriction inherits
  monotonicity, idempotence, and any WNU or Siggers identity of $f$.
\end{lem}
\begin{proof}
  $f$ preserves the order-convex $T$ (Lemma~\ref{lem:polchar}); an
  order-convex $C\subseteq T$ is order-convex in $\B^m$ (if $a\le c\le
  b$ with $a,b\in C\subseteq T$ then $c\in T$ by convexity of $T$, then
  $c\in C$), so $f|_T$ preserves it.  Identities restrict pointwise.
\end{proof}

Thus non-existence of a WNU or Siggers polymorphism on $\B^3$ propagates
to all $m\ge3$.

\section{The base case on $\B^3$}

\begin{prop}\label{prop:base}
  $\Pol(L_3)$ contains no ternary WNU operation.
\end{prop}
\begin{proof}[Proof (a twelve-cell contradiction)]
  Suppose $f$ is one.  For an antichain pair $\{x,y\}$,
  $\conv\{x,y\}=\{x,y\}$, so $m(x,y)\in\{x,y\}$; and $f(110,101,011)\in
  \{110,101,011\}$ (the coatoms are a $3$-antichain).  Consider the ten
  values
  \[\textstyle
  \begin{array}{llll}
  A=m(010,101), & B=m(100,011), & C=m(101,010), & D=m(101,011),\\
  E=m(101,110), & F=m(001,110), & G=f(110,101,011), & H=m(110,101),\\
  I=m(110,001), & J=m(110,011), & &
  \end{array}\]
  with domains the order-convex hulls
  ($A,C\in\{010,101\}$, $B\in\{011,100\}$, $D\in\{011,101\}$,
  $E,H\in\{101,110\}$, $F,I\in\{001,110\}$, $J\in\{011,110\}$,
  $G\in\{011,101,110\}$), and the monotonicities
  \[
    A\le G,H;\quad B\le D,G,J;\quad C\le D,E;\quad F\le E,G;\quad I\le H,J.
  \]
  As in Proposition~\ref{prop:nosiggers}, each is witnessed by two
  representatives comparable coordinatewise (the value is $f$ of any cell
  in its orbit, by the WNU identities):
  \begin{center}\small
  \renewcommand{\arraystretch}{1.05}\setlength{\tabcolsep}{4pt}
  \begin{tabular}{ccc@{\quad}ccc}
  \hline
  ineq. & rep.\ $\le$ & rep. & ineq. & rep.\ $\le$ & rep.\\
  \hline
  $A\le G$ & $(010,101,010)$ & $(110,101,011)$ & $C\le E$ & $(010,101,101)$ & $(110,101,101)$\\
  $A\le H$ & $(101,010,010)$ & $(101,110,110)$ & $F\le E$ & $(110,001,001)$ & $(110,101,101)$\\
  $B\le D$ & $(011,100,100)$ & $(011,101,101)$ & $F\le G$ & $(110,001,001)$ & $(110,101,011)$\\
  $B\le G$ & $(100,100,011)$ & $(110,101,011)$ & $I\le H$ & $(001,110,110)$ & $(101,110,110)$\\
  $B\le J$ & $(011,100,100)$ & $(011,110,110)$ & $I\le J$ & $(001,110,110)$ & $(011,110,110)$\\
  $C\le D$ & $(010,101,101)$ & $(011,101,101)$ & & &\\
  \hline
  \end{tabular}
  \end{center}
  If $B=011$ then $D=011$, $C=010$, $E=110$, $F=110$, $G=011$, and $F\le
  G$ fails ($110\not\le011$).  If $B=100$ then $D=101$, $C=101$,
  $E=101$, $F=001$, $G=101$, $A=101$, $H=101$, $I=001$, $J=110$, and
  $I\le J$ fails ($001\not\le110$).  Both cases are impossible.  (The
  twelve cells are a minimal unsatisfiable core located by machine; the
  deduction is by hand.)
\end{proof}

\begin{prop}\label{prop:nosiggers}
  $\Pol(L_3)$ contains no Siggers operation.
\end{prop}
\begin{proof}[Proof (an eight-variable contradiction)]
  Suppose $s\in\Pol(L_3)$ is a $4$-ary Siggers operation.  The Siggers
  identity $s(a,r,e,a)=s(r,a,r,e)$ groups the cells of $s$ into orbits;
  order-convex preservation puts each value in the convex hull of its
  arguments, and monotonicity relates comparable cells.  Consider the
  eight orbit-values
  \[
    P=s(110,001,001,110),\ Q=s(101,010,010,101),\ R=s(100,011,011,100),
  \]
  \[
    T=s(101,011,011,101),\ U=s(110,011,011,110),\ V=s(100,100,010,100),
  \]
  \[
    W=s(101,110,011,101),\ X=s(101,110,101,101),
  \]
  with domains the convex hulls
  \[
    P\in\{001,110\},\ Q\in\{010,101\},\ R\in\{011,100\},\ T\in\{011,101\},
  \]
  \[
    U\in\{011,110\},\ V\in\{010,100\},\ W\in\{011,101,110\},\ X\in\{101,110\},
  \]
  and the monotonicities
  \[
    P\le U,\ P\le W,\ P\le X,\ Q\le T,\ Q\le W,\ R\le T,\ R\le U,\ V\le W,\ V\le X.
  \]
  Each of these is witnessed by two orbit representatives that are
  comparable coordinatewise; since $s$ is constant on each orbit, the
  named value equals $s$ of the displayed representative, and
  monotonicity of $s$ gives the inequality.  The witnesses are:
  \begin{center}\small
  \renewcommand{\arraystretch}{1.05}
  \begin{tabular}{ccc}
  \hline
  inequality & representative $\le$ & representative\\
  \hline
  $P\le U$ & $s(110,001,001,110)$ & $s(110,011,011,110)$\\
  $P\le W$ & $s(110,001,110,001)$ & $s(110,101,110,011)$\\
  $P\le X$ & $s(110,001,001,110)$ & $s(110,101,101,110)$\\
  $Q\le T$ & $s(101,010,010,101)$ & $s(101,011,011,101)$\\
  $Q\le W$ & $s(101,010,010,101)$ & $s(101,110,011,101)$\\
  $R\le T$ & $s(100,011,011,100)$ & $s(101,011,011,101)$\\
  $R\le U$ & $s(100,011,011,100)$ & $s(110,011,011,110)$\\
  $V\le W$ & $s(100,100,010,100)$ & $s(101,110,011,101)$\\
  $V\le X$ & $s(100,100,010,100)$ & $s(110,101,110,101)$\\
  \hline
  \end{tabular}
  \end{center}
  (in each row the left argument tuple is $\le$ the right one in each of
  the four coordinates).  Split on $W$.  If $W=011$ then $P\le W$ forces $P=001$ and $V\le W$
  forces $V=010$; then $P\le X$ gives $X=101$, and $V\le X$ reads
  $010\le101$, false.  If $W=101$ then $P=001$, $Q=101$ (from $Q\le W$),
  $V=100$; $Q\le T$ gives $T=101$, $R\le T$ gives $R=100$, $R\le U$ gives
  $U=110$, and $P\le U$ reads $001\le110$, false.  If $W=110$ then
  $P=110$ and $Q=010$; $Q\le T$ gives $T=011$, $R\le T$ gives $R=011$,
  $R\le U$ gives $U=011$, and $P\le U$ reads $110\le011$, false.  All
  three cases are impossible.  (These eight orbits are a minimal
  unsatisfiable core located by machine; the deduction is by hand.  As a
  check the same search on $\B^2$ returns a Siggers operation.)
\end{proof}

\section{The dichotomy}

We record first the tractable side, for a single generator with
median-closed fibres.

\begin{prop}[median-closed fibres, after \cite{PaperW}]\label{prop:sw2}
  Let $z\colon\B^m\to\B^n$ be monotone with every fibre $z^{-1}(q)$
  closed under the ternary median $\mathrm{med}(x,y,z)=(x\wedge y)\vee
  (y\wedge z)\vee(z\wedge x)$.  Then the realizability problem along $z$
  carries the median as a majority polymorphism.  By the Baker--Pixley
  theorem \cite{BakerPixley} the solution relation of any instance is then
  determined by its binary projections; in the terminology of constraint
  satisfaction, a majority polymorphism yields \emph{strict width two}
  \cite{BartoKozik}: running the standard $(2,3)$-minimality
  (path-consistency) procedure on an instance produces a nonempty reduced
  instance exactly when the instance is solvable, decidable in polynomial
  time.  In particular the problem is in $\mathrm P$.
\end{prop}
\begin{proof}
  Applied coordinatewise the median preserves each fibre (they are
  median-closed) and the order $\{(x,y):x\le y\}$ (it is monotone), and
  satisfies the near-unanimity identities $\mathrm{med}(x,x,y)=
  \mathrm{med}(x,y,x)=\mathrm{med}(y,x,x)=x$; a ternary majority
  polymorphism gives strict width two \cite{BakerPixley,BartoKozik}.
\end{proof}

The contrast is total.

\begin{thm}\label{thm:cx}
  For $m\le2$ the order-convex language $L_m$ has bounded width (the
  median is a majority polymorphism) and $\mathrm{CSP}(L_m)\in\mathrm P$.
  For every $m\ge3$: $L_m$ has \emph{unbounded width}
  (Prop.~\ref{prop:base}, Lemma~\ref{lem:restrict}), $\Pol(L_m)$ is
  \emph{not Taylor} (Prop.~\ref{prop:nosiggers}, Lemma~\ref{lem:restrict}),
  and $\mathrm{CSP}(L_m)$ is $\mathrm{NP}$-complete.
\end{thm}
\begin{proof}
  The language $L_m$ contains every singleton $\{a\}=[a,a]$ (an
  order-convex set), so every retract of $L_m$ fixes each element:
  $L_m$ is a core and $\Pol(L_m)$ is idempotent.  The dichotomy
  \cite{Bulatov,Zhuk} in its idempotent formulation therefore applies
  directly, with no core reduction needed: $\mathrm{CSP}(L_m)$ is in
  $\mathrm P$ if $\Pol(L_m)$ is Taylor and $\mathrm{NP}$-complete
  otherwise.  For $m\le2$ the median is a Taylor (indeed majority)
  polymorphism: on $\B^{\le2}$ the coordinatewise median of any three
  vertices is one of the three ($\B^{\le1}$ is a chain, where this is
  clear; for $\B^2$, if the median differed from all three,
  each input would disagree with it in at least one of the two
  coordinates, so by the pigeonhole principle two inputs would disagree
  with it in the \emph{same} coordinate; but in that coordinate the median
  bit is the majority of the three input bits, which cannot differ from
  two of them---a contradiction.)  Being one of its arguments, the median
  preserves every unary relation, in particular every order-convex one,
  and it is monotone and idempotent---a majority polymorphism of $L_m$ by
  Lemma~\ref{lem:polchar}.  For $m\ge3$,
  by Lemma~\ref{lem:restrict} a ternary WNU
  or a Siggers operation of $\Pol(L_m)$ would restrict to one of
  $\Pol(L_3)$, which Propositions~\ref{prop:base} and~\ref{prop:nosiggers}
  exclude; so $\Pol(L_m)$ has no ternary WNU (unbounded width, by
  \cite{BartoKozik,KKVW}) and no Siggers operation (not Taylor, by
  \cite{Siggers,KMM}), and $\mathrm{CSP}(L_m)$ is $\mathrm{NP}$-complete.
\end{proof}

\begin{cor}[median-closure is the boundary for the canonical polymorphism]\label{cor:boundary}
  Median-closure of the fibres is the boundary for the \emph{canonical}
  coordinatewise median majority polymorphism of the realizability engine:
  \begin{enumerate}
  \item if a generator $z$ has median-closed fibres, realizability along
    $z$ has strict width two (Proposition~\ref{prop:sw2}) and is in
    $\mathrm P$;
  \item once a fibre is order-convex but not median-closed, the median is
    no longer a polymorphism, and the ambient order-convex constraint
    language $L_{m'}$ ($m'\ge3$)---the finite template with all
    order-convex relations---has unbounded width and is
    $\mathrm{NP}$-complete (Theorem~\ref{thm:cx}).
  \end{enumerate}
\end{cor}

Two cautions keep the statement honest.  First, part~(2) is a property of
the \emph{template} $L_{m'}$, not of any single generator: realizability
along a fixed $z$ uses only $z$'s own fibres and was polynomial in every
computed instance, so transferring the $\mathrm{NP}$-hardness to a uniform
cubical realizability problem would require a further reduction that we do
not carry out.  (The image $\im(z)$ need not be a sublattice---a monotone
$\B^3\to\B^3$ can have an $N_5$ image---so no distributivity is assumed.)  Second, the passage
from strict width two to a \emph{skeletally coskeletal} atom needs the
local-consistency step (three-dimensional realizability implies
$(2,3)$-consistency); Proposition~\ref{prop:medred} reduces it, for
median-fibre atoms, to a sharp pair-extension property that we verify
computationally but do not prove, so the coskeletal finiteness stays
conjectural (Part~\ref{part:a})---what is unconditional here is strict
width two itself.
What is unconditional is the algebraic dichotomy of
Theorem~\ref{thm:cx}, whose threshold is median-closure: Figure~\ref{fig:cube}
shows the elementary obstruction, the median of the three atoms escaping
their hull.

\begin{figure}[ht]
\centering
\begin{tikzpicture}[x=1.5cm,y=1.1cm,every node/.style={font=\footnotesize},
   dot/.style={circle,fill,inner sep=1.2pt},
   atom/.style={circle,draw,fill=black!12,inner sep=1.5pt}]
  \node[dot,label=below:$000$](b) at (0,0){};
  \node[atom,label=below left:$100$](a1) at (-1,1){};
  \node[atom,label=below:$010$](a2) at (0,1){};
  \node[atom,label=below right:$001$](a3) at (1,1){};
  \node[dot,label=above left:$110$](c1) at (-1,2){};
  \node[dot,label=above:$101$](c2) at (0,2){};
  \node[dot,label=above right:$011$](c3) at (1,2){};
  \node[dot,label=above:$111$](t) at (0,3){};
  \foreach \a/\c in {a1/c1,a1/c2,a2/c1,a2/c3,a3/c2,a3/c3}\draw[gray!55](\a)--(\c);
  \foreach \a in {a1,a2,a3}\draw[gray!55](b)--(\a);
  \foreach \c in {c1,c2,c3}\draw[gray!55](\c)--(t);
  \draw[-{Latex[length=2mm]},red!70!black,thick,dashed](a2) to[bend right=18] node[left,pos=.62]{$\mathrm{med}$}(b);
\end{tikzpicture}
\caption{The three atoms of $\B^3$ (shaded) form an antichain equal to
its own hull, but their median is the bottom $000$: no majority
operation compatible with the order-convex sets survives.}
\label{fig:cube}
\end{figure}
